\documentclass[12pt]{article}  
\newcommand{\say}[1]{``#1''}  
\newcommand{\nsay}[1]{`#1'}  
\usepackage{endnotes}  
\newcommand{\B}{\backslash{}}  
\renewcommand{\,}{\textsuperscript{,}}  
\usepackage{setspace}   
\usepackage{tipa}  
\usepackage{hyperref}  
\begin{document}  
\doublespacing  
\section{\href{reflection-2.html}{On Reflection}}  
First Published: 2026 July 5

\section{Draft 2: 5 July 2026}

Three years ago, at the end of this month, I wrote a \href{reflection}{post} about reflection.  
In that, I focused mostly on the chronology of my life as it pertained to writing and what the point of it was.  
A lot has changed since then; a lot has changed in the past months and days.  
So, let's take some time and reflect.

Yesterday I read a book that helped me understand something that I have struggled to articulate to many people.  
The book is called \say{The First 20 Hours}\footnote{already on the list of future posts as a book review}, and it explicitly calls out in the intro the idea that it takes ten thousand hours to master something.  
Rarely is my, or the author's, goal to master something, however.  
Instead, my goal is competency or at least mediocrity.  
The author and I agree that this is a much faster process, though he goes so far as to call it attainable in twenty hours.

He does an interesting job contrasting practice, learning, and skill acquisition, which I'm sure I'll touch on in the review.  
It helped me mostly in that it reminded me that I do not have to try to be the best in the world at everything I do; I can and want to and probably should seek instead to be able to do many things at least acceptably well.

I've also been thinking about what I want out of my life lately, especially because it's hard to be long distance from my partner.  
That's a situation we both want to resolve, but that involves uprooting my entire life.\footnote{just realized there's no section on social in the reflection. Adding now. Why social? because yesterday the bartenders were sad to hear I'm planning to uproot my life and I talked with the person who's acted as the sound and light manager for the bar for the past 13 years. I can't imagine myself in thirteen years, and he was lovely to talk to.}  
Before I tear the roots I've grown out of the ground, I want to know where they go.  
I'm also moving apartments, for the first time since moving for graduate school.  
I've lived here for six years, and that's not nothing, in the span of life.

I just realized I haven't really stopped to consider that this apartment where I'm currently writing is the place I've lived the second longest ever.  
For the next decade, at least, it will remain that.  
If I do not outlive my mother\footnote{though I do, of course, hope to and plan to}, this is where I will have lived for 10 percent of my life.  
That's something to say goodbye to, and a home to mourn for sure.  
Anyways, back to the reflection.

Why am I here?\\  
What is my calling?\\  
How do I get there?\\

These questions are ones that I consider often, and they are questions I do not have an answer to.

But, nearly six thousand words of writing later, I do also find that I'm out of the mood to reflect.  
Maybe writing by hand would help, but I doubt it.  
So, I'll choose to treat this time of reflection as what it was.  
I'm grateful I had the morning to reflect\footnote{coincidentally, it's currently 1159, so wow. (just turned 1200)}, and I'm grateful that I've had these past few weeks to reflect.

I am grateful for those in my life, past and present.\\  
I am grateful that I never, for a moment, doubted that I could come to my mother with anything to receive the help I need.\\  
I am grateful that I still know that there's always a place I can return if I need.

I am grateful.

I am still filled with grief about my mother's passing; I likely always will be.\\  
I am sad that I am losing the place I spent the past six years and that I'm also losing the time it takes me to move. This is the last place my mother will ever help me move into.\\  
I am unhappy that I do not live near my partner. I want to build a life with her, and our distance makes that harder.\\  
I am sad that my friends from graduate school have and are moving elsewhere, even as I am so excited and glad that they are building and continuing their lives. I am especially grateful for those who have kept me in their lives.\\  
I am not satisfied with my current work. It is not where I am meant to be for life. That's something that I knew going into the job, but it's still a shame.\\  
I'm nervous to go back to work tomorrow; I worry that people will be mad or disappointed in me for being injured, and I worry that I'll be unable to get the work done that needs to be finished.

I have many negative emotions, but they do not define me.

I am defined by the boundaries which create me.  
Those boundaries come primarily in the interactions I have with others and the bonds we made; I am where these bonds intersect.  
I am not the person I was at the start of this musing, and I am the person that I was when I laughed at my very first breath.

I find that this sort of reflective mood is not the most conducive to me doing anything, but I also find that I tend to drown it out.  
I'm going to take some time to see what it feels like to just sit with the feeling; I'll lie down and just breathe for a bit.  
After that, I'll work on removing the traces of my life in this home.\footnote{sorry that the very ending here got more practical again, but that's where my mind is at. This footnote is because the app I use wanted another five words}

\section{Draft 1: 5 July 2026}

I don't entirely know how I'm feeling right now.  
A few hours ago, maybe even an hour ago, I would have said solely positive.  
But, I realize that I had not given myself space to pause at all up to that point in the day.  
I woke up, rushed out (with an audiobook on) to Pilates, then left Pilates, left voicemails, turned back on the audiobook, and only turned it off to shower and shave, then write.  
At no point today was there a moment of just sitting and being.

I said that I would let myself take a nap after having some caffeine, because I do know that caffeine also helps when feeling sleepy.  
I drank the caffeine, and then I lied down\footnote{I think that lay/laid/lain is the one with an object and lie/lied is the one with out} for about ten minutes.  
In that time, I did some thinking, though about what I no longer fully remember.  
I did some resting, I imagine.  
And, more than anything, I had my headphones off\footnote{I tend to write with my headphones on because they're noise canceling}.

That didn't dim the mood.  
But, I do think that it made the mood at least a little more introspective.  
Reflective, one might say.

It's a good thing that this was already the plan for today's blog post!  
The last time I had a post with the URL \say{reflection} was about \href{reflection}{three years ago}.  
That post is not quite negative, though it does end on a negative note.  
It's focused much more on a chronology of my life and considerations for what I'm doing writing.  
It's strange to realize that it's been ages since I wrote in my little notebook at all, and so my penmanship hasn't gone through the same forced change it otherwise would have.

I've reintroduced lower-case letters to my writing, though.  
I think that came because I was going into pure\footnote{read: intentionally using cursive letters, rather than picking up the pen} cursive for a while, and then I wrote a letter to a child.\footnote{who, in retrospect, still couldn't read at all, nor should I have expected that}  
Children, as far as I'm to know, learn cursive well after print, and I wanted to be considerate of that.

I still write my letters strangely.  
My capital A is still closer to a triangle, and my lower case a is now effectively a 2.  
I'd been trying to avoid having hard reversals in the letters I write.  
When I started working on my left hand's hand\footnote{that's probably a redundant hand, but I'm not sure}, I noticed that cursive is full of these backtracking notes.  
It's strange, then, that it's considered the running script; I find that my running form doesn't have many hard stops and starts.  
Still, penmanship is something I'd like to get back into; I no longer get compliments about the way my writing looks, so something is clearly degrading.

I thought about the point of what I write during the last reflection.  
That's no longer a huge consideration for me.  
I said that this blog might be a nice resource for my future children; I neglected to consider that it's also a great place for a future partner to learn about me and how I've changed through the years/what's been on my mind.\footnote{in retrospect, I'm no longer even certain that I told my previous significant other (hmmm is there a reason that I use SO for one and partner for the other? probably) this blog existed. I'd actually be pretty certain I didn't tell her}

The penultimate paragraph of the post goes: \begin{quote}

When I look back at myself this time next year, what will I wish that I
spent more time on? I have to imagine that it will be similar things
that I wish past me had spent time on now. I wish that I would have
practiced art, though I can never really explain why. I wish that I
would be better at organizing. And, despite the fact that one of my most
common compliments is that I am good at keeping in touch, I wish I did better at that.  
\end{quote}

I didn't end up blogging a year later.  
The reflection post came at the end of July.  
I wrote a post \href{reflection-2024a}{July first} and \href{writers-slump}{July second}, and then not again until November.  
Even just looking at that first post when I came back is still hard, which is something that's going to be kept with me forever.

As a child, I think that at times it confused me why my dad still so clearly carried the weight of his own father's passing.  
After all, his father had last known him as a high schooler.  
At this point in my life\footnote{wow, three A sentences in a row, that's interesting}, I do not think of myself as particularly connected to even the college part of me, let alone the high school part of me.  
And yet, that if anything makes the loss of my own mother hit all the harder; she never got the chance to see me like this.\footnote{or, depending on our views of afterlife, I never got to have her see me like this, as the newly minted Dr. Rebelsky}  
The weight of that is something that I'm going to feel on many days; growing up I knew that my side of the aisle at the wedding would not be filled with first cousins\footnote{what most people tend to just call cousins, in my experience}, aunts, uncles, and grandparents.  
At this point, I at least have an aunt and an uncle, even if I don't have a mother I can walk down the aisle with.  
She will not be there to see her grandchildren born or her son's partner follow in her footsteps.

My goal with this post was not to wax melancholic about the passing of my mother.  
Then again, my goal with almost nothing these days is to consider her passing; that might be part of the issue.

When my grandmothers both died, I was in middle school.  
I kept myself from sinking into sorrow the best way that I knew how; I listened to music as soon as I woke up and until I went to sleep.  
It went so far that I tried to use a song looping as an inverse alarm clock, thinking that when the song stopped playing I might wake to its absence.  
(For the record, that did work, though it was not a restful night).  
This lasted until my mother and I went to the Boundary Waters for a canoeing trip.

There are no electronics on that trip, or at least there weren't when I went on it.  
The books that were brought were explicitly intended as firestarters; when we were done with half the book, those pages would be used to start the flames for our nightly fire.\footnote{for those feeling bad, they were mass-market paperbacks, we didn't burn the plastic covers, and they were recent enough books that I am certain there existed other copies}  
We talked to the other people on the expedition, we played games, and of course, we worked out.  
It's not nothing to portage a canoe and two weeks of supplies every day for two weeks.

But, there was also a lot of time for silence.  
I cried a lot that fortnight, and it was very healthy for me.  
When I returned home, I still listened to music, but it was no longer to drown out my thoughts.

These days, I can ignore the lyrics in music.  
I can, in fact, almost ignore all of music, if I try.  
It still acts as something which prevents my thoughts from fully going, though.  
Likewise, the headphones I wear, although noise-cancelling, produce an inverse sound that works in a similar way, allowing me to block out part of what I don't want to consider.  
If I'm going to be reflective right now, then, I should take off the headphones.\footnote{wow I immediately feel uncomfortable. Why didn't I bring the weighted blanket with me for safety?}

Three years ago, I talked about how I was being ramblier than usual in the post.  
Looking at my posts lately, they're all at least as rambly as that one was.\footnote{my dictionary doesn't like rambly, probably rambling would be the ideal term, but I would say rambly, and I want this to be an accurate reflection of me, at least in part}  
That's something that's interesting to note, and I'm sure I'd notice many other things if I was to continue reading through old posts.  
I do especially appreciate my note about three sentence paragraphs, which I do often find myself using.

Looking at the Draft 0, I see a few things I wanted to talk about.  
First is something that will be irrelevant by the time any of my readers can see this post; the abridged version is that the back-end for my blog was written in 2018.  
The code it runs on is now deprecated, and so I need to update something to make the posts get pushed to the site.

The second comes from a conversation with my partner.

This morning, my Pilates instructor was a former graduate student in the same department as me.  
She Mastered\footnote{since Master is the degree title and is capital, I'm capitalizing the verb as well} out, and she now works as a Pilates instructor and professional dancer.  
Lately I've been thinking that I want to use my Ph.D. in my long-term work, in a way that my current job doesn't seem equipped to allow.  
But, thinking about how much joy it brought me to see a chemist doing something else, I know also that I don't want to do only chemistry.  
No, that's not totally it.

I've known since the start of my Ph.D. and before that I am not someone who has a single interest.  
Or, at least, the interest I have and had are not those which map onto a career that becomes my life, much as I admire people like Erd{\H o}s.  
I especially admire those thinkers who, by drifting through different fields, are able to impact multiple fields.

I think about this in the context of a podcast I listened to recently about Natural Law.  
The podcast was about how the modern Catholic apologetics around the theory tend to ignore a lot of what is foundational to both the conception in Aristotle and Aquinas's Catholicization of the idea.  
I'm not going to rehash the ideas in the podcast here,\footnote{though you could watch it on YT \href{https://www.youtube.com/watch?v=apF41LTl6Lw}{here}} but something that the author mentioned is really meaningful to me; the great thinkers in Natural Law did so while also practicing other pieces of human flourishing: poetics and poetry.  
Few, if any, of the leading apologists who weaponize\footnote{because when wielding to dismiss others, it's a weapon} Natural Law are also practicing poets or work in poetics.\footnote{which I understand is the study of poetry as craft and art and philosophy}

I mentioned on the call with my partner that I think that I want to work in chemistry the way that an 18th century person might have.  
Not, as joked by her, in a shed with no rules, but rather as something integrated into the quest for knowledge.  
Three years ago, I thought about how the writing I do serves me, and I asked if it's bringing me closer to G-d.  
In a way, it's a similar question here; how does me using my Ph.D. bring the world to G-d?\footnote{I can tell that some of my philosophy is straying from the Roman orthodoxy, in part with the fact that I think that theosis (wow it took me a while to find it), the becoming like G-d, is doable and what we should strive for, as the Eastern Church does. Also, like the Jews of my ancestry, I take the idea that my service is to the world and of the world, rather than personal salvation. that's something to think about. (adding to the post list)}

Or, at least, it feels like a similar question.  
Chemistry does not provide the tools to answer every question.  
Where chemistry provides the tools, it's incredibly powerful.  
I do love rotational spectroscopy, and I do love astrochemistry.  
I don't know if I see a world where my fulfillment comes in absence of both of those.\footnote{though, admittedly, that's quite possibly just me not being able to envision right now}

How, exactly, chemistry impacts my life will depend so much on the life I end up living.  
If all proceeds according to current plans, I'll be married with children, likely serving in more of a domestic role.  
I know from one of my mathematician friends that some of what I consider baseline ingrained action is seen by the outside world as chemistry; they messaged me all excited about the kitchen chemistry they did to make cheap Pedialyte replacement.  
I know that I want to instill in my children and other children the idea that the world is understandable.  
I know that I feel a deep need to make knowledge more accessible.

I have been wanting to get into philosophy for more than a year now.  
When the idea of finishing my dissertation was still more an amorphous idea than something with a deadline, I really wanted to spend time looking at the ways that philosophy can interact with science.  
No, that's not right.  
I wanted to look at the ways that the Academy I exist within fundamentally has its sets of philosophies and how we can improve them for the sake of improving knowledge, or at least how we can expose them for the uninitiated.

I still want to do that, and I still think about the fact that it's only in the oldest quantum mechanics textbook I found that there's an intuitive explanation for bonding.\footnote{I do not recall the book, but I'm sure I'll find it. It mentions that we know many things \textit{a priori}, such as that there is an ideal bond length somewhere between nothing (because we are not an infinitely dense mass) and infinite (because we are not infinitely separated particles). Adding to the posts on teaching (astro)chem doc}

In making the astrochem doc and adding the text, I realized that I have a folder of old blog posts.\footnote{admittedly, not well organized, and not well-filled}  
I have a post I was going to reflect on here, \say{On Grief}.  
However, I have not posted this anywhere.  
That's interesting, and I did start to go through to see what else I have.

That's not the project for now, though.  
Nor is reading my post about grief or the idea of turning 60, though I am interested in those as well.\footnote{and I'm adding cleaning this site to the list of tasks to do in the future}

Where was I?

eh, I've rambled off track here enough that I'm going to move on to a new draft.   
Maybe Draft 2 will be coherent?

\section{Draft 0.5: 5 July 2026}

I'm planning to write this draft before I do my daily reflections, because I generally want to spend some time thinking about where I am.

It's always fun to see the ways that my posts can interact in conversation with each other; three years ago at the end of the month, I also \href{reflection.html}{wrote a blog post about generally being reflective}.  
I'm sure that the main reason I'm feeling reflective this morning is that I recognize that today is the last day before I return to my normal behavior.\footnote{read: working at my corporate office job}  
It also helps that I woke up early, went to a lovely Pilates class, and am generally feeling good today.\footnote{which I'll attribute in part to eating something last night}

So, let's quickly think about what I wanted to reflect about?

No, I've gotten distracted from that.\footnote{before the second line of Draft 0, my partner called and we had a lovely 20ish minute call. I was already kind of distracted then, so I am even less sure what I wanted to write then}

Let's do the daily reflection and then go back to this to write the major reflection.\footnote{or, most likely, drink some more water and eat the bagel I bought and maybe take a nap}

Hmm upon finishing the daily reflection, that is a two thousand word daily reflection. That's like (45 minutes?) of writing, which is maybe more time than I should be spending. I guess that a few hundred of the words are template, and that's ok. Still, I don't know how conducive this will be to my every day experience. I guess album club won't have a full draft each day. Also like, I think that all the questions are good questions. It still only took me an hour to do, including the breaks and whatnot, which is an hour that I'm happy to spend thinking about life.

One question will have to be whether we can do this writing in the morning? Might be good to start the day with that kind of intentionality.

Regardless, I'm going to take a break before the next draft and order my agenda for the day, then go and close the many HumbleBundle tabs I have open right now.

\section{Draft 0: Some random rambling I did on 5 July 2026}

Ugh the blog from yesterday failed to post... well, I guess that we'll have to do some git training today

reflection plan today: what do i mean by chemistry like 1800s? Like integrated with other things

Bring up the whole "natural law, aquinas was also a poet" , so was aristotle

book from yesterday, easy to get competent. I already have/had most of the skills.  
One of the skills he learned was coding for his website... i should maintain my own site

\langle this text was removed because it's not something I want publicly shared right now \rangle

\section{Upcoming Posts}  
\begin{itemize}

\item Wonder  
\item Curse of knowledge, esp re. Music Theory\footnote{stolen from a prior list}.

Most of a draft done, though rambly.  
\item Overview of a course on astrochem

I think that I have writing somewhere about it.  
\item RebelFit

A rambly draft done, and some more work which tells me that I was not doing anything wrong in the past, approximation just takes time.  
\item Why aspiring writers should practice typing\footnote{and this should also like tie in the thing that a lot of writers are told to do which is directly write good poetry or prose down to get it in the hand. This also has the effect of practicing hand writing. I think that there's also the whole like \say{when we want to be better at running, we also stretch }}

\item At some point I'd like to reflect on my hesitance to delete email, which I think comes from a place of digital hoarding. That's going to be a fun one for sure.

\item Moving (homes)

\item Algorithmic constraints in art I create\footnote{see footnote in \href{reflection-june-2026}{this post's second daily reflection about written word, poetry}}

\item Singing in a low-voice choir, especially in context of my other music

\item Listening to content at speeds (J speed versus single speed)\footnote{see footnote in \href{reflection-2}{this post's daily reflection about written word, audio book} for context}

\item Book Review(s):\footnote{also added today as a list so I remember}

\begin{itemize}

\item First 20 Hours

\item Dungeon Crawler Carl

\end{itemize}

\item Ideas of Salvation\footnote{from the First Draft of \href{reflection-2}{this post}, footnote where I discuss changes to my theology.}, in particular, how mine have changed over time.

\end{itemize}

\section{Daily Reflection: 5 July 2026}

\begin{itemize}

\item All in all, how goes it?

Honestly, really good. I'm feeling great right now, though kind of hungry and thirsty, which I'm about to address with food and drink before starting the actual line by line items.\footnote{I feel like I'm a little too live-streaming this post. I did feel like I lost an hour, but in retrospect, probably not}

\item Creative Considerations:

\begin{itemize}

\item Gut and / or vibes check, how is creative life going?

About as well as yesterday, which means that this set of reflections may not be useful as an everyday thing.

\item Written Word:

\begin{itemize}

\item How's your relationship with the written word in general?

Really good!

\item How goes blogging? What's feeling like the barrier?

Honestly right now the barrier is nonexistent.  
Big issue I'm having right now is that my blog is written in a now deprecated version of something so it can't deploy.

\item How goes reading?

Really good! I read two books last night, a romance novel that a friend lent me\footnote{I think that after the bartender invites me to after-work drinks I get to call them a friend} and a book about how to learn competency in skills quickly.  
The competency book I skimmed more towards the end, because I already know how to do ukulele and don't want to develop an interest in windsurfing.\footnote{because I know that I'd want to do it so much and it's way too dangerous for a boi like me right now}

\item How goes writing poetry again?

Generally decently.  
As we saw in \href{reflection-june-2026}{yesterday's post,} I am able to write an (admittedly bad) poem in the form I want relatively quickly.  
That's the end of what I did, though, so that's not the ideal for me, especially long term.  
I think maybe getting back into a sonnet a night could be nice?

\item We're counting audio book here, anything good?

Still finishing Bedlam Bride, 42 minutes and 47 seconds left.  
OH, I should reflect about listening to audiobooks at single speed.  
Added to the list.

\item Upon reflection, how's the written word? Anything missing in the four\footnote{fixed today} questions?

Generally decent! I think that I'd like to get back into reading about writing again.  
I should add a line about the web serial again. That's something that I want to return to, so having a daily reminder.  
No updates there, so adding that to the list for tomorrow's reflection.\footnote{look at me, optimistically assuming I'll be blogging tomorrow}

\end{itemize}

\item Music

\begin{itemize}

\item In general, how goes music?

Generally ok! I still don't think that I'm having the strongest relationship with music, but that's also sometimes the ebb and flow of life.

\item How goes recording?

None

\item How goes writing?

None, but that's something I want to do today (and so I'm adding it to today's fungenda)\footnote{fungenda is a neologism I learned from my partner. It's a portmanteau of fun and agenda}

\item How goes playing?

Didn't do any, which is sad.  
The book about learning things\footnote{wow that's coming up a bunch} did end with a chapter on learning ukulele, so I did feel some connection to music then, if nothing else.  
My favorite bar also has live music on Fridays\footnote{I think that's the plural}, and that's something that could be fun to consider for some Friday night if I get back into playing a whole lot more.\footnote{and, presumably, have a backing band}  
I do miss playing in front of others.

\item How goes listening?

Well, I have listened to the full album for album club this week. Let's take this space to write the reflection:

Album of the week was the 2022 remaster of The Who\footnote{I'm choosing to capitalize the}'s \textit{Who's Next}

I always forget how \say{Baba O'Riley}'s intro does nothing to really connect to the vibe of the rest of the song, even if the arpeggiated organ sound repeats through the whole song. Didn't listen closely enough to see if it's actually a completely static through line, but it felt like it was! I like the song.

As I already mentioned to the group\footnote{this is getting sent to them directly because I don't entirely know how to figure out how to fix the blog}, \say{Getting in Tune} feels more modern than the rest of the album.  
There's a podcast/YouTube channel I've recently started watching (Actually Sounds Like with Jim Lill),\footnote{who also makes a series on where the sounds in recordings actually come from that's really cool, if only because I find it really impressive when non-classically trained scientists do really rigorous science} and his first episode was on the Beatles.  
In that, he mentioned how remastering can do a lot to change the sound of a song, and I heard what he meant when playing two versions of one of the Beatles' songs.  
Not sure if that's why?

\say{The Song is Over} feels like a death ballad, in some ways.\footnote{I'd forgotten which sound it was, which is on me}

\say{Behind Blue Eyes} was probably my favorite, in part because it seemed like a song I could recreate myself. It's primarily guitar and vocals in harmony, which is a sound I really like.

In total, though, the album felt really like one sound through the whole time, and I didn't give it the time to listen to each song on repeat until I really felt the feel of what made each unique.  
As with most older albums, what really strikes me is how much time they let things develop and play.  
\say{Baba O'Riley} doesn't introduce anything but the arpeggio until more than thirty seconds in, when the crashing piano overlays it.  
They give the piano three and a bit repetitions (another fifteen seconds) before introducing the drums.  
This five minute long song doesn't start vocals until one minute and six seconds in.  
I can never imagine a modern song doing that, because people would click out well before then.

But, it's something that I really admire.

\item Upon reflection, how goes music? Anything missed in those four questions?

I think that reflecting on music did a lot to help me feel more connected. I guess there's something to be said for \say{did I listen to the Jim Lill stuff?} but that's probably not something that I really care too deeply about.

\end{itemize}

\item Embroidery Project?

\begin{itemize}

\item In general, what's the vibe of the project right now?

Good! I wrapped my hoop yesterday which was nice.

\item How goes progress on the design? What's in the way?

No progress since last post. I spent some time last night while in bed thinking about the base layer, the color gradient.  
I'm not sure what I want it to look like, and the issue with doing two corners of warm and cold colors is the diagonals.  
Issue with doing a single point in the middle is the same.  
Issue with having a wall is that I then have a full line of a single color which I don't really want.

I'll probably just need to play around with it for a while.

\item How goes progress on the stitching? What's in the way?

At this point, I feel pretty confident about most of it. I'm going to be doing four total strands (two doubled).

Hmm, I did really really like the 14 strand with 6 that I did before... Ugh Ok I guess that I'll see whether the pain I was having was just because of the material having a bunch of junk under it.

I wrapped my hoop yesterday though! That was really good of me, and I'm proud of myself.

\item Upon reflection, how goes embroidery? Anything missed in those two questions?

In general, slowing down a touch, but that's not something I'm unhappy with. I think it's good for me to take things as they come. I do have some questions about how I'll carry the design around with me and whatnot, but that's something I can consider later.  
I have a nice big embroidery frame that would likely hold the entirety of the project, but I'll need to figure out how to attach the project to the frame, and it's a little too big to just casually carry around, which is an issue of its own.

Eh, need to finish the design first anyways.

\item Project Progress\footnote{new today, so all action items new today. In general I'll plan to add action items with footnotes as needed}

\begin{itemize}

\item Base layer design:

Still in ideation mode

\item Weave layer design:

Have demonstrated crossing with three stitches wide and four stitch gap at widest point. Looks nice, shows enough of the background that I'm happy with it, probably.  
Will need to consider in context of the top layer

\item Top layer design:

Have the knot made\footnote{thanks isometric drawing guide!}, and have a braid pattern that works for the six colors going straight up. Still need:

\begin{itemize}

\item Figure out how to do turns  
\item Figure out if the pattern I have for the braid works with angles that aren't straight up  
\item Figure out how to do the braid going sideways  
\item Fill in the knot with the braid and make sure it looks ok  
\item Decide if I'm going to mark every single stitch (probably), and if so, start marking in the braid stitches

\end{itemize}

\item Full design: Figure out how many colors of thread I'm going to buy and how much of each

\item Full design: Figure out how it all looks together

\item Full design: plan out the action list for making it happen

\item Stitching: do

\end{itemize}

\end{itemize}

\item Outreach:

\begin{itemize}

\item How's outreach going?

Nonexistent.

\item How goes prep for the rest of your life?

Nonexistent. Planning to write up a thing I can send to libraries and whatnot about what I can do, though.

\item How goes prep / submission of workshops?

Added to the fungenda.

\item Any talks coming?

Sadly no :(.

\item Upon reflection, how goes outreach? Anything missed in those three questions?

Nonexistent right now, but I'm getting back into thinking about it, at least.

\end{itemize}

\item Upon reflection, how is creation going?

About what I thought! That's kinda nice, not going to lie.

\end{itemize}

\item Health Considerations:

\begin{itemize}

\item Gut and / or vibes check, how is health going?

I will always laugh at me putting \say{gut check} for my health, because something something the gut is part of health.  
In general, not sleeping the best, but feeling generally healthier?

\item Working out?

I did a cardio Pilates class this morning, which was really fun! It absolutely set me up for the day, though I am hitting the midday sleepies.\footnote{read: it's 10 AM and I'm feeling my eyes get heavy. Not sure if it's a lack of food or water or sleep or caffeine, but I'll probably lie down for a bit after this daily reflection}

\item Food?

Still not as great as it could be. Yesterday I ate a cream puff over six or so hours\footnote{admittedly, a BIG cream puff}, an order of fries, and about half a large pizza.\footnote{the pizza felt like not much food at all, but in retrospect, half a pizza is probably a meal, especially because it's like 1800 calories, if I'm doing my math correctly. Ok I feel better about it now}

This morning before Pilates I had a cinnamon sugar stick. Since Pilates, I've had a mozzarella and hard salami on rye bagel sandwich.

\item Water?

At the bar yesterday, I had what I think is averaging out to at least a glass of water an hour. That's not a great sign, given how little water I've been drinking lately. Probably good for me though! Right now I'm sipping on my bottle as I write.

\item Sleep?

I went to bed much earlier than I have been last night!\footnote{about 9:15}

I woke up much earlier than I have been this morning!\footnote{out of bed by 6:50}

I still toss and turn at night, but the weighted blanket is absolutely helping. When I'm up in the middle of the night, I still don't necessarily know what I need to do to fall back asleep. I'm hoping that the combination of Sonnet Time and me working out this morning will have helped.

\item Upon reflection, how is it going?

Decent! I think that the place that's going best right now relative to before is the mental load. I feel excited to do today, and that's not something I've been able to say a lot of days. It probably helped that I met a chemist today at Pilates.

\end{itemize}

\item Practical Considerations:

\begin{itemize}

\item Gut and / or vibes check, how are the practical things going?

I don't think as well as they could be, but not bad!

\item How goes packing old home?

Ehhhhhh. I've taken over a load of dirty clothes.\footnote{new place has a washing machine IN UNIT!!!!}  
I would like to spend time today figuring out how to rent a moving van so I can take my table (and probably bed, though I'm hopeful that will be easier to move).

\item How could you be working on your LDR better?

I feel like I've been monopolizing our conversations more than is the ideal, because I would love to hear more about my partner's life and we have such little time to catch up.\footnote{read: an hour is never enough, and ten isn't either}

\item How goes setup of new home?

I haven't really done anything for that. I've continued to think about how I want it to be arranged, and I should really look into getting some sort of furniture for friends to come over and sit on?

\item The work you're being paid for?

Starts tomorrow! I have nerves for sure.

\item Upon reflection, how go practical matters?

Could be better. Big action item today I think will be being more intentional about how much I'm speaking during the call time.

\end{itemize}

\item Upon total reflection, how goes it?

Good! About what I thought, but that's also not something that I think is really that bad.

\end{itemize}

\end{document}